\scanpage{43}
\margin{42}
Pour pallier à cet inconvénient on peut songer à ``multiplier'' chaque $D\in\Sigma$, en autant d'exemplaires qu'il y a de sommets, de sorte que \deleted{chaque} pour un $(D,s)$ donné (``élément de contact'' de $\underline\Sigma$), il y aura exactement deux \uline{arcs de} \deleted{spécialisation} \uline{générisation} \inserted{\uline{exceptionnelle}} correspondants sur $\Delta_D$, qui n'empiètent pas. Si le sommet $s$ est simple, alors les deux arcs de générisation exceptionnels sont les deux côtés d'une structure de bigône topologique, Sinon, ces arcs, plus les \inserted{deux} arcs de générisation \inserted{(ordinaire)} correspondants à $s$, forment une structure de polygone topologique carré sur $\Delta_D$. \deleted{En} Dans les deux cas, il n'y a pas d'arc de' inertie --- qu'à cela ne tienne !

On a \inserted{tout} ce qu'il faut pour compléter une surface compacte, en recollant à la précédente des bigones ou carrés. \deleted{Mais il faudrait déterminer} \inserted{En \unsure{terme} de graphes,} \deleted{\unsure{les}} on a remplacé le graphe (orienté) précédent, par celui déduit en remplaçant chaque sommet ``exceptionnel'' par l'ensemble des \deleted{\unsure{diamètres}}\inserted{\uline{diagonales}} du polygone local
